\honest{Tolerate Tolerance}{Eliezer Yudkowsky}{2009}

People who set out to be rationalists are likely to be more annoyed than most by flaws in reasoning. It ``doesn't strictly follow,'' I grant, but ``a lot of us probably'' have the annoyance set high. So if we want to get anything done together, we must tolerate other people's tolerance.

My main case is Ben Goertzel, who has something nice to say about everyone. He even complimented the ideas of ``the most legendary of all AI crackpots,'' who then, apparently, added a link to the compliment to his email signatures. \nb{The FAQ the post links quotes the compliment, and adds that Goertzel later conceded critics' points and, in 2004, said he had no definite opinion of the ideas.} I have come to see this as one of Goertzel's strengths: he is nice to people others might ignore, me included, and now and then it pays off for him. If I mark him down for it, I am insisting that he dislike everyone I dislike before I will work with him, which is unrealistic when different people are annoyed by different things in different amounts. I admit it is hard to remember this when he is nice to ``so many idiots.''

Cooperation needs some punishment of defectors. But punishing people for refusing to punish is ``a far more dangerous idiom that can lock an equilibrium in place even if it's harmful to everyone involved.'' \nb{Boyd and Richerson's model of 1992 found exactly this: punishing non-punishers can make ``any individually costly behavior'' stable, ``whether or not it creates a group benefit.''} I remind myself of this when Robin Hanson points out a flaw in some academic trope and modestly confesses he could be wrong (he is not), or Michael Vassar sees promise in someone I dismissed in thirty seconds.

I do get annoyed when others give too much credit, and I suspect some of my fellow rationalists do too; it may even be universal, with ``an obvious evolutionary rationale'' that I do not spell out, which would make it an unpleasant and dangerous adaptation.

I am not a fan of ``tolerance,'' and I do not believe in being ``intolerant of intolerance,'' ``as some inconsistently hold.'' I give no argument that it is inconsistent. \nb{Karl Popper's ``paradox of tolerance'' (1945) argues that ``Unlimited tolerance must lead to the disappearance of tolerance,'' and so presents intolerance of intolerance as a limit tolerance needs.} But I will tolerate people more tolerant than I am, and judge them only for their ``own un-borrowed mistakes.''

If a group waits for you to hate the right enemies, you may be in the wrong group. Just do not demand that all your friends be as intolerant of that group as you are. Forgive them if some suggest that Group X was not so awful after all.
